% =====================================================================
% Introduction.
%
% The Pingora framework this backend integrates with, why this document
% exists, the artifact types and the relationships between them, and the
% flow-of-refinement organisation the rest of the chapters follow.
% =====================================================================

\section{The Pingora Framework}

Pingora is an open-source Rust framework, developed and open-sourced by
Cloudflare, for building fast, reliable, and programmable networked
systems.\footnote{\href{https://github.com/cloudflare/pingora}{cloudflare/pingora on GitHub}}\textsuperscript{,}\footnote{\href{https://blog.cloudflare.com/pingora-open-source/}{Cloudflare blog: ``Open sourcing Pingora''}}%
\textsuperscript{,}\footnote{\href{https://developers.cloudflare.com/changelog/post/2026-05-04-pingora-powers-cache/}{Cloudflare changelog, 4~May 2026}}
It was created as a modern, memory-safe alternative to C/C++-based proxies and has been
battle-tested in production at Cloudflare, where it serves more than
40~million Internet requests per second. Pingora is built entirely in Rust on
top of an asynchronous runtime, and provides the building blocks for HTTP/1
and HTTP/2 proxies, gRPC and WebSocket proxying, load balancing, graceful
reloads, and rich observability.

The framework is organised as a collection of crates, each addressing a
distinct concern. The most relevant to this specification are:

\begin{itemize}
    \item \texttt{pingora-core} --- transport protocols and the core traits
          that define server and service behaviour.
    \item \texttt{pingora-proxy} --- the HTTP proxy state machine and the
          \texttt{ProxyHttp} extension points used to customise proxy logic.
    \item \texttt{pingora-cache} --- the HTTP caching layer, including the
          storage backend abstraction that this work targets.
\end{itemize}

\section{Purpose of This Document}

This document specifies the requirements for a \emph{file-system-backed}
implementation of the \texttt{pingora-cache} \texttt{Storage} trait. By
persisting cached objects on disk rather than solely in memory, the cache can
retain entries across process restarts and scale well beyond the limits of
main memory, while remaining a drop-in backend behind the existing
\texttt{HttpCache} state machine. The driving goals are:

\begin{itemize}
    \item \textbf{Persistent storage} --- use the file system as a durable
          medium for cached HTTP objects so that the cache survives restarts
          and grows beyond RAM capacity.
    \item \textbf{Standards compliance} --- preserve the RFC~9111 and HTTP
          caching semantics already enforced by \texttt{pingora-cache}, so
          persistence does not compromise correctness.
    \item \textbf{Efficient use of space} --- store and reclaim cached objects
          economically, supporting eviction and avoiding unnecessary disk
          overhead.
\end{itemize}

\section{The Refinement Model}

\subsection*{Artifacts and Relationships}\label{sec:artifacts}

The artifact types and the relationships between artifacts take inspiration
from \sysml, the Object Management Group's Systems Modeling
Language,%
\footnote{\href{https://www.omg.org/}{\omg}}%
\footnote{\href{https://www.omg.org/sysml/}{\sysml}}%
\footnote{P.~Roques,\href{https://re-magazine.ireb.org/articles/modeling-requirements-with-sysml}{``Modeling
Requirements with \sysml''}, \emph{Requirements Engineering Magazine}, IREB,
April~2015, section ``Requirement Relationships''.}%
without aiming to follow it
strictly.

Every artifact has a type, given by the prefix of its ID, and every type carries
a \emph{stereotype}: a label written between guillemets, such as
<<stakeholder requirement>> or <<design constraint>>, that names the type's
category. Table~\ref{tab:sysml-requirement-stereotypes} lists every artifact
type with its prefix and stereotype. Two types
can share a stereotype --- functional designs (\textsf{FDS}) and functional
requirements (\textsf{FRE}) are both <<functional requirement>> --- so it is the
prefix that tells an artifact's kind. The diagrams draw each artifact as a card
headed by its stereotype, and each relationship as an arrow labelled with the
relationship's name, pointing from the artifact that depends to the one it
depends on.

\includereqstereotypes

Each relationship has one job in this model:
\begin{itemize}
  \item \emph{Satisfies} (<<satisfy>>): a decision or a design satisfies the
    requirements it answers.
  \item \emph{Derives} (<<deriveReqt>>): a requirement derives from the
    artifact it comes from. Usually that is a broader requirement it breaks
    down, as a functional requirement breaks down a stakeholder need. It can
    also be a decision: choosing a mechanism creates constraints of its own
    that no stakeholder asked for, and those derive from the decision that
    created them (for example, choosing a plain file system,
    \reqref{ADR-PERSISTENCE.MECHANISM-01}, is what makes the number of entries
    in one directory matter, so \reqref{REQ-EFFICIENCY.DIRSCALE-01} derives
    from that decision). A requirement can have both kinds of parent.
  \item \emph{Refines} (<<refine>>): a decision or a design refines the broader one it works
    out in more detail. A design normally refines the decision it carries out.
    It may instead refine another design, when it builds on terms that design
    defines and no decision of its own lies between the two (for example, the
    design that logs and counts metadata that outgrew its record,
    \reqref{FDS-OBSERVABILITY.HEADROOM-01}, refines the design that defines when
    that happens, \reqref{FDS-STORAGE.METAUPDATE-01}). The link to the decision is not
    lost: the other design refines its own decision in turn.
  \item \emph{Verifies} (<<verify>>): a test case verifies the requirement,
    decision, or design it checks.
  \item \emph{Traces} (<<trace>>): records any other dependency that none of the
    above fits.
  \item \emph{Conflicts with} (<<conflicts with>>): marks
    two artifacts that cannot both be met as stated; the artifact that declares
    the conflict also states how it is to be resolved.
\end{itemize}
Tied together this way, the artifacts form a chain that can be read forwards
(goal to test) or traced
backwards (test to goal); the Traceability Analysis chapter renders it
explicitly. Table~\ref{tab:sysml-relationship-types} summarises the
relationships and the direction of their arrows.

% A float rather than a page-breaking xltabular: the table is short, and
% splitting it strands its last rows on the next page. [tb], like the
% generated stereotype table, so neither is set alone on a float page.
\begin{table}[tb]
\centering\small\renewcommand{\arraystretch}{1.3}%
\caption{Relationships between artifacts}\label{tab:sysml-relationship-types}
\begin{tabularx}{\linewidth}{|l|>{\raggedright\arraybackslash}p{0.3\linewidth}|X|}
\hline
\textbf{Relationship} & \textbf{Arrow (from $\rightarrow$ to)} & \textbf{Meaning} \\
\hline
\textit{<<satisfy>>} & Decision, design, or implementation $\rightarrow$ requirement & Fulfils the requirement \\
\hline
\textit{<<deriveReqt>>} & Derived requirement $\rightarrow$ its source & Follows from the source --- a broader requirement it breaks down, or a decision whose choice created it \\
\hline
\textit{<<refine>>} & Detailed artifact $\rightarrow$ broader artifact & Works the broader artifact out in more detail \\
\hline
\textit{<<verify>>} & Test case $\rightarrow$ artifact & Checks that the artifact holds \\
\hline
\textit{<<trace>>} & Any $\rightarrow$ any & A dependency no more specific relationship describes \\
\hline
\textit{<<conflicts with>>} & Either way & The two cannot both be met as stated; the declaring artifact states how the conflict is to be resolved \\
\hline
\end{tabularx}
\end{table}

\subsection*{Flow of Refinement}

Rather than grouping its contents by artifact type --- all requirements in one
chapter, all architecture decisions in another --- this document is organised to
follow the \emph{flow of refinement}: the order in which the artifacts are
discovered as a need is progressively elaborated from an abstract goal down to a
verifiable design. Each chapter is one stage in that flow, and any stage may
contain artifacts of \emph{any} type; what a chapter groups is a \emph{degree of
refinement}, not a kind of artifact. The stages are preceded by a \emph{System
Context} chapter that documents, non-normatively, the \texttt{pingora-cache}
behaviour the artifacts cite as their basis.

The stages, and the artifact types that typically appear in each, are:
\begin{enumerate}
  \item \textbf{High-Level Requirements} (\textsf{HLR}, and use cases
        \textsf{UC}) --- the stakeholder needs the cache must meet, stated without
        commitment to a solution.
  \item \textbf{High-Level Refinement} --- the first elaboration of those needs:
        the foundational architecture decision that shapes the whole solution
        (an \textsf{ADR}, \emph{Architecture Decision Record}, at its broadest
        scope). A decision, here or later, \emph{generates} requirements of its
        own --- stated as base requirements (\textsf{REQ}) when left unbranded,
        or as functional requirements (\textsf{FRE}) when they state, in
        verifiable terms, \emph{what} the backend must do; either kind may
        instead derive directly from the stakeholder need itself. The later,
        more detailed decisions must satisfy whichever of these the earlier
        ones generated.
  \item \textbf{Detailed Refinement} --- the decisions that follow from the
        high-level one: further architecture decisions (\textsf{ADR}) at a
        narrower scope, down to a single mechanism.
  \item \textbf{Structural and Functional Design} --- concrete designs realising
        those decisions: structural designs (\textsf{SDS}, e.g.\ the on-disk
        layout) and functional designs (\textsf{FDS}, e.g.\ sequence diagrams).
  \item \textbf{Test Case Specifications} (\textsf{TC}) --- the checks that verify
        the requirements and designs of the preceding stages.
\end{enumerate}
Table~\ref{tab:sysml-requirement-stereotypes} in Section~\ref{sec:artifacts}
lists the artifact types in this same order.

The stages order the artifacts; two further axes tell them apart. The first is
\emph{altitude}: an \textsf{ADR} can sit anywhere from the foundational choice
of storage medium down to a single mechanism.
The second is \emph{kind}: a \emph{requirement} (\textsf{HLR}, \textsf{REQ},
\textsf{FRE}) states a need without choosing a solution; a \emph{decision}
(\textsf{ADR}) argues a choice among alternatives
--- recorded as context, alternatives, and consequences --- and closes design
freedom; a \emph{design} specifies the chosen result --- \textsf{SDS} its
static shape, \textsf{FDS} its dynamic behaviour --- and no longer decides
anything; and a \emph{test case} (\textsf{TC}) verifies that what was
required, decided, or designed actually holds.

Two tests keep the kinds apart in practice. The \emph{deletion test}: if
removing an artifact would force a trade-off to be argued again, it is a
decision; if the trade-off would stay settled but the implementer would lose
the blueprint, it is a design. The \emph{change-trigger test}: a decision is
reopened when its forces change --- a new constraint, a new alternative, a
shifted trade-off --- while a design is revised when the decision stands but
its realisation needs adjustment.

\subsection*{Domains}\label{sec:domains}

Every ID also names a \emph{domain}: the subject its artifact belongs to,
written between the type prefix and the first dot.
\reqref{HLR-PERSISTENCE.DURABLE-01} belongs to \textsf{PERSISTENCE}, and
\reqref{FDS-STORAGE.PUBLISH-01} to \textsf{STORAGE}. A domain cuts across the
stages: it holds the requirements of one concern together with the
decisions, designs and test cases that answer them, wherever in the document
they appear.

Figure~\ref{fig:domain-overview} shows the domains, how many artifacts of
each type each one holds, and how they depend on one another. An arrow counts
the relationships from the artifacts of one domain to those of another, by
kind, in the colours the diagrams give each relationship; \emph{trace} is
left out. Each domain is placed by the stage of refinement its artifacts
belong to, on average and test cases aside, so the domains of requirements
stand above those of the decisions and designs that answer them; an arrow
points down only where a domain of an earlier stage depends on one of a
later stage. Every artifact closes with a link to its domain in this figure;
its own diagram and its list of relationships lead on to its neighbours.

\begin{figure}[htbp]
  \centering
  \makebox[\textwidth][c]{\reqoverviewdiagram}
  \caption{Domain overview}\label{fig:domain-overview}
\end{figure}
